\documentclass[12pt]{article}

\usepackage[a4paper, margin=1in]{geometry}
\usepackage{amsmath, amssymb}
\usepackage{setspace}
\usepackage{titlesec}
\usepackage{enumitem}
\usepackage{hyperref}

\setstretch{1.15}

\titleformat{\section}{\large\bfseries}{\thesection}{1em}{}
\titleformat{\subsection}{\normalsize\bfseries}{\thesubsection}{1em}{}

\title{\textbf{CSE400 -- Fundamentals of Probability in Computing}}
\author{Lecture 4: Joint Probability and Conditional Probability}
\date{Instructor: Dhaval Patel, PhD \\ January 15, 2026}

\begin{document}

\maketitle
\hrule
\vspace{0.5cm}

%------------------------------------------------

\section{Course Context and Motivation}

\subsection{Why Learn CSE400?}
Daily life conversations involve uncertainty.

\subsection{Engineering Applications}
\begin{itemize}
\item Speech Recognition
\item Radar Systems
\item Communication Networks
\end{itemize}

\subsubsection*{Speech Recognition System}
Speech templates for words (Hello, Yes, No, Bye) are compared to recognize spoken input across different speakers.

\subsubsection*{Radar System}
Radar systems detect objects by transmitting signals and analyzing reflections.

\subsubsection*{Communication Network}
Information travels from a \textbf{source} to a \textbf{destination} through interconnected nodes.

%------------------------------------------------

\section{Lecture Outline}

\subsection*{Introduction to Probability Theory}
\begin{itemize}
\item Experiments, Sample Space, and Events
\item Axioms of Probability
\item Corollaries and Propositions
\item Assigning Probability: Classical and Relative Frequency
\end{itemize}

\subsection*{Joint Probability}
\begin{itemize}
\item Motivation, notation, and concepts
\item Card deck example
\item Costume party example
\end{itemize}

\subsection*{Conditional Probability}
\begin{itemize}
\item Motivation and notation
\item Cards without replacement
\item Game of poker
\item The missing key
\end{itemize}

\textit{Note: Provided slides cover the Introduction to Probability Theory section.}

%------------------------------------------------

\section{Introduction to Probability Theory}

\subsection{Experiments, Outcomes, and Events}

\subsubsection*{Definition: Experiment ($E$)}
A procedure we perform that produces some result.

\textbf{Example:} Tossing a coin five times ($E_5$).

\subsubsection*{Definition: Outcome ($\xi$)}
A possible result of an experiment.

Example outcome:
\[
\xi_1 = HHTHT
\]

\subsubsection*{Definition: Event}
A set of outcomes of an experiment.

Example event $C$ for $E_5$:
\[
C = \{\text{all outcomes consisting of an even number of heads}\}
\]

%------------------------------------------------

\subsection{Sample Space}

\subsubsection*{Definition: Sample Space ($S$)}
The collection of all possible distinct outcomes of an experiment.

\textbf{Properties}
\begin{itemize}
\item \textbf{Mutually Exclusive:} Only one outcome occurs at a time.
\item \textbf{Collectively Exhaustive:} All possible outcomes are included.
\end{itemize}

\textbf{Notes}
\begin{itemize}
\item $S$ is the universal set of outcomes.
\item Sample space may be:
\begin{itemize}
\item Discrete
\item Countably infinite
\item Continuous
\end{itemize}
\end{itemize}

\subsubsection*{Examples of Sample Spaces}
\begin{enumerate}
\item Flipping a fair coin once
\item Rolling a cubical die
\item Rolling two dice
\item Flipping a coin until a tails occurs
\item Random number generator over interval $[0,1]$
\end{enumerate}

%------------------------------------------------

\section{Probability and Axioms}

\subsection{Probability}
Probability is a measure of the likelihood of events, or a function that assigns a numerical value to an event representing its likelihood.

\subsection{Axioms of Probability}

\textbf{Axiom 1 (Bounds)}
\[
0 \leq Pr(A) \leq 1
\]

\textbf{Axiom 2 (Certainty)}
If $S$ is the sample space:
\[
Pr(S) = 1
\]

\textbf{Axiom 3 (Additivity)}
If $A \cap B = \emptyset$:
\[
Pr(A \cup B) = Pr(A) + Pr(B)
\]

For infinitely many mutually exclusive events $A_i$:

\[
Pr\left(\bigcup_{i=1}^{\infty} A_i\right) =
\sum_{i=1}^{\infty} Pr(A_i)
\]

%------------------------------------------------

\section{Corollary from Probability Axioms}

\textbf{Corollary 2.1}

For a finite number of mutually exclusive events $A_i$, $i=1,2,\dots,M$:

\[
Pr\left(\bigcup_{i=1}^{M} A_i\right)
= \sum_{i=1}^{M} Pr(A_i)
\]

\textbf{Notes}
\begin{itemize}
\item Equivalent to Axiom 3 for finite sample spaces.
\item Axiom 3 is required for infinite sample spaces.
\end{itemize}

%------------------------------------------------

\section{Propositions from Probability Axioms}

\subsection{Complement Rule}
\[
Pr(A^c) = 1 - Pr(A)
\]

\subsection{Monotonicity}
If $A \subset B$, then:
\[
Pr(A) \leq Pr(B)
\]

%------------------------------------------------

\section*{Exam Summary}

\begin{itemize}
\item Experiment: procedure producing a result.
\item Outcome: possible result.
\item Event: set of outcomes.
\item Sample space: all possible outcomes.
\item Outcomes are mutually exclusive and collectively exhaustive.
\item Probability follows three axioms.
\item Complement rule and subset relation follow from axioms.
\item Finite additivity simplifies calculations.
\end{itemize}

\end{document}
